\newcommand{\hmwkTitle}{Homework 1}
\newcommand{\hmwkDueDate}{Tuesday, January 27, 2026}
\newcommand{\hmwkHeaderTitle}{ICCS206 Discrete Mathematic: Homework 1}
\documentclass{article}

\usepackage{fancyhdr}
\usepackage{amsmath}
\allowdisplaybreaks
\usepackage{amsthm}
\usepackage{amsfonts}
\usepackage{tikz}
\usepackage[plain]{algorithm}
\usepackage{algpseudocode}
\usepackage{enumitem}
\usepackage{graphicx}
\usepackage{hyperref}

\usetikzlibrary{automata,positioning}

%
% Basic Document Settings
%

\topmargin=-0.45in
\evensidemargin=0in
\oddsidemargin=0in
\textwidth=6.5in
\textheight=9.0in
\headsep=0.25in

\linespread{1.1}

\pagestyle{fancy}
\lhead{\hmwkAuthorName}
\chead{} % empty
\rhead{\hmwkHeaderTitle}
\cfoot{\thepage}

\renewcommand\headrulewidth{0.4pt}
\renewcommand\footrulewidth{0pt}

\setlength\parindent{0pt}

%
% Homework Details
%

\providecommand{\hmwkTitle}{Homework}
\providecommand{\hmwkDueDate}{TBD}
\providecommand{\hmwkClass}{ICCS206 Discrete Mathematic}
\providecommand{\hmwkClassTime}{Section\ 1}
\providecommand{\hmwkClassInstructor}{Pakawut Jiradilok}
\providecommand{\hmwkAuthorName}{\textbf{Jiraroj Wiruchpongsanon (6781617)}}
\providecommand{\hmwkHeaderTitle}{\hmwkClass:\ \hmwkTitle}

%
% Title Page
%

\title{
    \vspace{2in}
    \textmd{\textbf{\hmwkClass:\ \hmwkTitle}}\\
    \normalsize\vspace{0.1in}\small{Due\ on\ \hmwkDueDate}\\
    \vspace{0.1in}\large{\textit{\hmwkClassInstructor\ \hmwkClassTime}}
    \vspace{3in}
}

\author{\hmwkAuthorName}
\date{}

\renewcommand{\part}[1]{\textbf{\large Part \Alph{partCounter}}\stepcounter{partCounter}\\}

%
% Helper Commands
%

\newcounter{partCounter}
\newcommand{\solution}{\textbf{\large Solution}}

\newtheorem*{theorem}{Theorem}

\theoremstyle{remark}
\newtheorem*{remark}{Remark}

\begin{document}

\maketitle
\newpage


\section*{Problem 1}
A number $x$ is a perfect square if and only if 
\[\exists m \in I \text{ such that } x = m^2\]
Prove that if $x$ and $y$ are both perfect squares, then $xy$ is also a perfect square.


\medskip
\solution

 \begin{proof}
  from the definition of a perfect square, given $x$ and $y$ are
  perfect square, we say
  \[
  x = n^2
  \]
  and
  \[
  y = m^2
  \]
  now, we can write $xy$ as
  \begin{align*}
  x \cdot y &= n^2 \cdot m^2\\
  % idk what's wrong with the spacing, but it doesn't look that bad
  &= (nm)^2 && \text{by the commutative and associative property
  of }\mathbb{R}
  \end{align*}
  thus $xy$ is also a perfect square
  \end{proof}

\vspace{1em}
\hrulefill


\section*{Problem 2}
If $n = ab$, then $\sqrt{a} \le \sqrt{n}$ or $\sqrt{b} \le \sqrt{n}$.

\medskip
\solution

\begin{proof}
  this one, we do contrapositive proof, let the $P$ be proposition $n
  = ab$, Q be proposition $\sqrt{a} \le \sqrt{n}$, R be proposition $
  \sqrt{b} \le \sqrt{n}$. The conjunction we're trying to proof can be
  written as
  \[
  P \to (Q \lor R)
  \]
    using contrapositive --- the logical property of $\to$
  \[
  P \to (Q \lor R) \equiv (\neg Q \land \neg R) \to \neg P
  \]
  which translate to
  \[
  \text{if }\sqrt{a} > \sqrt{n}\text{ and }\sqrt{b} > \sqrt{n},
  \text{ then }n \ne ab.
  \]
  since $\sqrt{a} > \sqrt{n}$ and $\sqrt{b} > \sqrt{n}$
  \[
  ab = (\sqrt{a}\sqrt{b})^2 > (\sqrt{n}\sqrt{n})^2 = n^2 \ne n
  \]
  therefore, $(\neg Q \land \neg R) \to \neg P$ is true, and thus $P \to (Q
  \lor R)$ is true
  \end{proof}

\vspace{1em}
\hrulefill


\section*{Problem 3}
Show that if $x$ is a rational number and $y$ is a rational number, then $x + y$ is a rational number.

\medskip
\solution

 \begin{proof}
  we define rational number as any number in the field (commonly $
  \mathbb{R}$ or $\mathbb{C}$) that can be written in a form of
  fraction of two integer with non-zero denominator --- so x and y can
  be written as
  \[
  x = \frac{x_n}{x_d}
  \]
  \[
  y = \frac{y_n}{y_d}
  \]
  now, if we sum up the two
  \begin{align*}
  x + y &= \frac{x_n}{x_d} + \frac{y_n}{y_d}\\
  &= \frac{x_n \cdot y_d + y_n \cdot x_d}{x_d \cdot y_d}
  \end{align*}
  with the closure of multiplication and addition property in $
  \mathbb{Z}$, each of $x_n \cdot y_d$, $y_n \cdot x_d$, and $x_d \cdot
  y_d$ stays an integer, and the numerator sum remains in $\mathbb{I}$.
  also, $x_d \ne 0$ and $y_d \ne 0$ imply $x_d \cdot y_d \ne 0$, so the
  fraction is valid.

  so, we can conclude that $\frac{x_n \cdot y_d + y_n \cdot x_d}{x_d
  \cdot y_d}$ is a fraction of integer, and thus x + y is a rational
  number
  \end{proof}

\vspace{1em}
\hrulefill


\section*{Problem 4}
Show that if $x$ is irrational then $1/x$ is also irrational.

\medskip
\solution

 \begin{proof}
  we will assume that x is rational number for the sake of
  contradiction, from the definition x can be written as
  \[
      x = \frac{a}{b} \text{ when } a,b \in \mathbb{I}
  \]
  then
  \begin{align*}
      \frac{1}{x} &= \frac{1}{\frac{a}{b}}\\
  &= \frac{b}{a} && \text{where $a$ and $b$ are still an integer}
  \end{align*}
  since $\frac{1}{x} = \frac{b}{a}$ --- a fraction of two integers;
  it implies that $1/x$ is a rational number, which contradicts the
  given that x is irrational, thus if x is irrational then $1/x$ is
  also irrational.
  \end{proof}

\vspace{1em}
\hrulefill


\section*{Problem 5}
Recall the card trick shown in class (\url{https://youtu.be/UmlZaONp-_I?t=3671}).
Suppose that we have a deck of $52$ cards (26 reds and 26 blacks). We draw two cards at a time until the deck is empty while counting the number of black pairs, red pairs, and different-color pairs. Show that the number of red pairs and the number of black pairs are always equal at the end.

\medskip
\solution

 \begin{proof}
  there are 3 scenario that could happened here
  \begin{itemize}
      \item a pair of red: notated with $R$
      \item a pair of black: notated with $B$
      \item mixed: notated with $M$
  \end{itemize}
  for set of red cards (26 in total), the set is a combination of
  \[
  2R + M = 26
  \]
  and for the set of black cards
  \[
  2B + M' = 26
  \]
  since a mixed pair must include both red and black card, $M = M'$,
  and we can solving the two equation gives $R = B$ which concludes
  that the number of red pairs and the number of black pairs are always
  equal at the end.
  \end{proof}

\vspace{1em}
\hrulefill


\section*{Problem 6}
Recall the absolute value function
\[
|x| = \begin{cases}
x & x \ge 0,\\
-x & x < 0.
\end{cases}
\]
Prove the statement
\[
\left|\frac{a}{b}\right| = \frac{|a|}{|b|}\qquad \forall a,b \in R,\ b \ne 0.
\]

\medskip
\solution

 \begin{proof}
  let us consider a, we'll have two cases:
  \begin{enumerate}
      \item $a \ge 0$
      \item $a < 0$
  \end{enumerate}
  consider case 1 first, there'll be two more sub-case if we consider
  b:
  \begin{enumerate}
      \item[$1.1$] $b \ge 0$ --- since both a and b are at least 0
      \[
          \frac{a}{b} \ge 0 \implies \left|\frac{a}{b}\right| = \frac{a}{b},\
  \text{and } |a| = a\ \text{and } |b| = b \implies \frac{|a|}{|b|} =
  \frac{a}{b}
      \]
      the theorem is satisfied
      \item[$1.2$] $b < 0$
      \[
      \frac{a}{b} \le 0 \implies \left|\frac{a}{b}\right| = -\frac{a}{b},\
  \text{and } |a| = a,\ |b| = -b \implies \frac{|a|}{|b|} = -\frac{a}
  {b}
      \]
      the theorem holds
  \end{enumerate}
  consider case 2, $a < 0$, again two sub-cases for b:
  \begin{enumerate}
      \item[$2.1$] $b \ge 0$
      \[
          \frac{a}{b} \le 0 \implies \left|\frac{a}{b}\right| = -\frac{a}{b},\
  \text{and } |a| = -a,\ |b| = b \implies \frac{|a|}{|b|} = -\frac{a}
  {b}
      \]
      the theorem is satisfied
      \item[$2.2$] $b < 0$
      \[
          \frac{a}{b} \ge 0 \implies \left|\frac{a}{b}\right| = \frac{a}{b},\
  \text{and } |a| = -a,\ |b| = -b \implies \frac{|a|}{|b|} = \frac{a}
  {b}
      \]
      the theorem holds
  \end{enumerate}
  since we have covered every case, we are done with the proof
  \end{proof}



\vspace{1em}
\hrulefill


\section*{Problem 7}
\begin{enumerate}[label=\arabic*.]
    \item State the contrapositive and prove: If $r$ is irrational, then $r^{1/5}$ is irrational.
    \item Is the following statement true? Prove or disprove: If $r$ is rational, then $r^{1/5}$ is rational.
    \item Is the following statement true? Prove or disprove: If $r^{1/5}$ is irrational, then $r$ is irrational.
\end{enumerate}

\medskip
\solution

\vspace{1em}
\hrulefill


\section*{Problem 8}
Consider the following game with $10{,}000$ players.
\begin{enumerate}[label=\arabic*.]
    \item Everyone starts with 1 million Baht.
    \item Each round shows a time-series graph of a stock to everyone.
    \item Each person may bet 1{,}000 Baht that the stock price will go up, go down, or stay the same exactly one hour from now, or they may opt out.
    \item After one hour the result is revealed and all bet money is split equally among those who guessed correctly (e.g., if 10 bet up, 20 bet down, and 5 bet same; if the stock goes up, each of the 10 winners receives $\frac{(10+20+5)\times 1{,}000}{10} = 3{,}500$ Baht).
    \item The game repeats from step 2.
\end{enumerate}
Your technical-analysis friend offers a program that reads the graph and tells everyone exactly what to bet, claiming everyone will make a profit after 20{,}000 rounds if everyone uses the program. Prove this claim cannot be true. (Hint: zero-sum.)

\medskip
\solution

\vspace{1em}
\hrulefill


\section*{Problem 9}
Consider a room of area $5\,\text{m}^2$ in which you can pick any shape, and you wish to place nine rugs each of area $1\,\text{m}^2$, also with shapes of your choice. Intouch demands that no pair of rugs may intersect with area at least $1/9\,\text{m}^2$. Prove that satisfying both Billy's and Intouch's requirements is impossible. (Hint: contradiction.)

\medskip
\solution

\vspace{1em}
\hrulefill


\section*{Problem 10}
Let $O$ be the set of odd integers. Prove that $\forall x \in O\ \exists p,q \in I$ such that $x = p^2 - q^2$; i.e., every odd integer can be written as the difference of two squares.

\medskip
\solution

\vspace{1em}
\hrulefill


\end{document}
